\section{Literature Review}

\subsection{2022 AML Closed-Vocabulary SSR}

The 2022 Imperial College AML Lab project demonstrated that a functional silent speech interface can be built using inexpensive surface EMG sensors and an Arduino microcontroller \citep{sun2022}. Their system used three MyoWare sensors placed on key facial muscles and applied classical preprocessing such as wavelet denoising, transient-state removal, and sliding-window feature extraction. A range of machine learning models were tested, from Random Forests, SVMs, and KNN to CNN architectures, with accuracies varying from 76 to 86 percent depending on the ML model. Upon testing new subjects, accuracy sharply dropped to 31.55\%.

\subsection{Foundational EMG-to-Speech Work}

The foundational benchmark for modern silent speech recognition was established by \citet{gaddy2021}, who developed an end-to-end system for converting facial EMG signals from eight electrodes into speech audio. Their approach predicts mel-frequency cepstral coefficients (MFCCs) from raw EMG and synthesises audio using a WaveNet vocoder. Two key architectural contributions drove their improvements: replacing hand-designed EMG features with learned convolutional features extracted directly from raw signals, and substituting LSTMs with a Transformer encoder using relative position embeddings. An auxiliary phoneme prediction loss provided additional supervision during training. To handle the temporal misalignment between silent and vocalised speech, the model uses dynamic time warping (DTW) to align silent EMG recordings against their vocalised counterparts during training. With approximately 40 million parameters trained on 19 hours of single-speaker data, the system achieved a word error rate (WER) of 42.2\% on open-vocabulary transcription, a substantial reduction from the previous 68.0\%. This work defined the dataset, electrode placement protocol, and evaluation methodology that most subsequent studies build upon.

\subsection{LLM-Augmented Silent Speech Recognition}

The current state of the art in unvoiced silent speech recognition on the Gaddy and Klein benchmark is the 2025 study \textit{``Can LLMs Understand Unvoiced Speech?''} \citep{mohapatra2025}. This work proposes a two-stage EMG-to-LLM pipeline in which an adaptor network maps high-frequency EMG into the input embedding space of a frozen large language model (LLaMA-2-7B and LLaMA-3.2-3B). The adaptor is composed of 1D convolutional layers for downsampling the high-rate EMG signal ($\leq$800 Hz), residual blocks for temporal feature extraction, BiLSTM layers for modelling sequential dynamics, and fully connected layers that project the extracted features into the embedding dimension of the target LLM. The resulting sequence of EMG embeddings is concatenated with text prompts and fed into the frozen LLM for autoregressive decoding.

Evaluated on the closed-vocabulary subset of the Gaddy and Klein dataset (67 words, approximately 26 minutes of data), the architecture achieves a WER of 0.49 with only 6 million trainable parameters, significantly outperforming application-specific EMG models with 54 million parameters that reach only 0.75 WER. The study further demonstrates that handcrafted EMG features improve LLM performance over raw signals, and that LLMs learn more effectively from audio than from EMG, highlighting the inherent difficulty of the unvoiced modality. An additional finding is that person identification from unvoiced EMG achieves 96\% accuracy, which motivates the development of personalised rather than speaker-independent models.

\subsection{Further Developments in Silent Speech Research}

Beyond the core works discussed above, several parallel developments have shaped the current landscape of silent speech research. \citet{benster2024} introduced MONA LISA, a cross-modal framework that aligns EMG and audio representations through contrastive learning objectives (crossCon and supTcon), performing temporal alignment in learned latent space rather than on raw signals. Combined with LLM-based rescoring of beam search outputs at inference time, the system achieves 12.2\% WER on silent EMG---the first to break the 15\% threshold for non-invasive silent speech---and 3.7\% on vocal EMG. In a different linguistic context, \citet{song2023} demonstrated that Transformer decoders substantially outperform LSTM decoders for syllable-level sequential decoding of Mandarin Chinese from 64-channel high-density sEMG, achieving 5.14\% CER and 96.37\% phrase accuracy. \citet{xie2025} further extended the field to Chinese with the first end-to-end silent EMG-to-text system, using a Transformer encoder with CTC decoding, auxiliary pinyin generation, and adversarial session classification to achieve 38.0\% CER across five recording sessions. Their ablation study confirms that facial electrode channels contribute significantly more than laryngeal channels in silent mode, consistent with the intuition that silent speech involves mouthing without vocal fold vibration.

On the practical side, \citet{inoue2025} address the problem of heterogeneous electrode configurations by proposing tokeniser strategies that allow a single HTNet+Conformer model to handle varying sensor setups, achieving 95.3\% word classification accuracy on a 64-word vocabulary and demonstrating cross-language transfer from Japanese to English. \citet{hwang2026} complement this with a systematic analysis of channel efficiency on the Gaddy dataset, showing that optimal reduced-channel subsets exploit complementary spatial coverage rather than simply retaining the individually strongest channels. Their channel dropout pre-training strategy, in which channels are randomly masked during full eight-channel training before fine-tuning on a target subset, consistently outperforms training from scratch in four-to-six channel configurations---a finding directly relevant to our project, which operates with fewer electrodes than the Gaddy benchmark. The emergence of self-supervised foundation models further advances the field: \citet{fasulo2026} introduce TinyMyo, a 3.6M-parameter Transformer encoder pre-trained via masked reconstruction that achieves competitive performance across gesture classification, kinematic regression, and silent speech tasks (33.54\% WER on speech synthesis) while being the first EMG foundation model deployed on an ultra-low-power microcontroller. Finally, the paradigm of projecting biosignal embeddings into frozen LLM token spaces extends beyond EMG: \citet{mishra2025} demonstrate with Thought2Text that EEG signals evoked by visual stimuli can be decoded into text using instruction-tuned LLMs fine-tuned through a CLIP-aligned multi-stage pipeline, validating the generality of the adaptor-based approach used by \citet{mohapatra2025} across different biosignal modalities.
